\documentclass{article}

\usepackage{arxiv}

\usepackage[utf8]{inputenc} % allow utf-8 input
\usepackage[T1]{fontenc}    % use 8-bit T1 fonts
\usepackage{hyperref}       % hyperlinks
\usepackage{url}            % simple URL typesetting
\usepackage{booktabs}       % professional-quality tables
\usepackage{amsfonts}       % blackboard math symbols
\usepackage{amsmath}
\usepackage{nicefrac}       % compact symbols for 1/2, etc.
\usepackage{microtype}      % microtypography
\usepackage{graphicx}
\graphicspath{ {./images/} }

\title{Density-Based Clustering and Contextual Graph Autoencoders for\\
Anomaly Detection in Industrial and Satellite Telemetry}

\author{
 Alessandro Bevilacqua \\
  Department of Computer Science \\
  University of Bari Aldo Moro \\
  Bari, Italy \\
  \texttt{a.bevilacqua11@studenti.uniba.it} \\
}

\begin{document}
\maketitle

\begin{abstract}
Multivariate telemetry from industrial plants and spacecraft comprises large, 
complex datasets, where anomaly detection represents one of the most critical and highly requested tasks.
This work reimplements DENCAST, a distributed density-based clustering algorithm built on locality-sensitive hashing, 
in PySpark, and
extends it into an anomaly detector, DENCAD, in which a graph autoencoder reconstructs
the centre of the contextual subgraph a new instance is attached to. DENCAD is compared with 
four other detectors on the same data and under the same protocol: 
the \emph{plain autoencoder}, the same network on one observation in isolation (graph-free); 
the \emph{graph autoencoder}, the same network over the $L$-hop neighbourhood the graph defines, without any clustering; 
\emph{DENCAST distance}, the distance to the nearest clustered object; 
and \emph{K-Means}, the distance to the nearest centroid.
All five produce a score and share one thresholding rule, 
a multiple of the training deviation whose factor is chosen on validation (sigma rule); 
the comparison itself rests on average precision and ROC-AUC, which are threshold-free, 
so they reflect the ranking each detector produces rather than where its threshold happens to fall.
From the results, it is quite clear that the autoencoder detectors are generally ahead of the distance-based ones, K-Means and DENCAST distance, even if not so much. 
Among the autoencoders no clear winner emerges, and the neighbourhood graph does not always justify its overhead. 
However, it is worth noting that DENCAD comes first on both ESA variants (normal and contaminated)
that count rare events as legitimate when considering AP as the metric.

\end{abstract}

\section{Introduction}

Industrial control systems and spacecraft emit continuous multivariate telemetry whose
faults are rare but costly. Detection is posed as a one-class problem: a model of normal behaviour is fitted on a period believed to be
clean and later instances are scored by how far they fall from it. Three properties make the setting challenging: 
the data is large --- the ESA mission used here holds seven million readings over fourteen years --- 
so the neighbourhood graph can be very expensive to compute; 
the detector must be \emph{inductive}, scoring an instance that arrives after training without refitting anything; 
and an anomaly is not always a property of one reading but often a trend across many.

DENCAST~\cite{corizzo2019dencast} addresses the first two challenges for clustering. It
approximates a neighbourhood graph with locality-sensitive hashing, identifies core
objects by their degree in the DBSCAN sense~\cite{ester1996dbscan} and propagates cluster
labels over the graph, all in Spark. Furthermore, a new object is assigned to an existing cluster
by similarity to a clustered one rather than by refitting.

The third challenge is addressed in the preprocessing stage of this work.
Readings are grouped into fixed, non-overlapping windows of a set length, and each window is reduced to summary statistics per feature:
mean, standard deviation, minimum and maximum for the numeric channels, and mode, entropy, distinct count and largest share for the categorical ones. 
The window is then the object that is hashed, clustered and scored, and it is labeled as positive if any of its constituent readings is positive (i.e. anomalous).

DENCAD constructs the contextual subgraph for a new window by retrieving its nearest clustered windows, 
and employs a graph autoencoder to reconstruct the target window from this context, using the reconstruction error as the anomaly score.

Our main contributions are a PySpark reimplementation of DENCAST starting from the original implementation written in Scala; DENCAD, an inductive detector
over contextual subgraphs using density-based clustering; and an evaluation protocol applying the same
thresholding rule and calibration procedure to five detectors.

\section{Related Work}

\paragraph{Distributed density-based clustering.} DBSCAN~\cite{ester1996dbscan} defines clusters by density, 
but its neighbourhood queries parallelise poorly.
DENCAST~\cite{corizzo2019dencast} replaces exact neighbourhoods with LSH candidates and
label propagation in order to adopt a design that fits Spark's shuffle model. We keep
its graph construction and clustering. Its hash family, random hyperplanes~\cite{charikar2002simhash}, is canonical for angular similarity
and invariant to the norm of the vector; since our features are individually standardised, that norm represents the aggregate deviation from normal operation---precisely the signal we want to preserve.
We therefore adopt the $p$-stable family of Datar et al.~\cite{datar2004pstable} (the family of E2LSH), whose collision probability decreases with Euclidean distance.

\paragraph{Graph autoencoders for anomaly detection.} DOMINANT~\cite{ding2019dominant}
reconstructs both attributes and structure of an attributed graph with a graph
convolutional encoder~\cite{kipf2017gcn} and ranks nodes by reconstruction error. It is
transductive: a new node requires recomputing the graph. Inductive
aggregation~\cite{hamilton2017graphsage} removes that constraint in principle; we obtain
it by construction, attaching the new window to a frozen graph and reading only its own
receptive field. Containing the cost of a full-graph convolution by partitioning is the
argument of Cluster-GCN~\cite{chiang2019clustergcn}; our contextual subgraphs exploit the
same economy, with partitions derived come from the density clustering rather than from a
graph partitioner.


\section{Materials}

Three public datasets are used, summarised in Table~\ref{tab:data}, where the ESA positive label can combine anomalies and rare events 
(i.e., rare events can be considered either legitimate or anomalous). 
SWaT is a water treatment testbed with labelled attack campaigns~\cite{mathur2016swat}; HAI~23.05 is a
hardware-in-the-loop industrial security dataset~\cite{shin2020hai}; ESA is Mission~1 of
the ESA Anomaly Dataset~\cite{kotowski2024esa}, comprising fourteen years of satellite telemetry
whose 200 annotations distinguish 118 anomalies, 78 \emph{rare events} (namely, unusual but legitimate behaviour), 
and 4 communication gaps. Splits are chronological, so validation and test follow the training period with no overlapping windows.

\begin{table}[t]
  \caption{Datasets after preprocessing and the windows they produce. \emph{Period} is the
  sampling interval and \emph{steps} the readings a complete window holds; $d$ is the number
  of window features. Positives and segments are counted on the test split, a segment being a run of consecutive positive windows.}
  \centering
  \small
  \setlength{\tabcolsep}{4pt}
  \begin{tabular}{lrrrrrrrrr}
    \toprule
    & \multicolumn{4}{c}{source} & \multicolumn{5}{c}{windows} \\
    \cmidrule(lr){2-5} \cmidrule(lr){6-10}
    & rows & sensors & period & steps & $d$ & train & valid. & test & pos./seg. \\
    \midrule
    SWaT &   946{,}719 & 52 &  1\,s &  60 & 124 &  8{,}280 & 3{,}749 & 3{,}750 & 728 / 13 \\
    HAI  &   802{,}800 & 87 &  1\,s &  30 & 220 & 19{,}080 & 3{,}840 & 3{,}840 & 169 / 19 \\
    ESA  & 7{,}364{,}160 & 76 & 60\,s & 132 & 216 & 34{,}576 & 8{,}509 & 8{,}629 & 937 / 24 \\
    \bottomrule
  \end{tabular}
  \label{tab:data}
\end{table}

\paragraph{Cleaning.} Rows with no reading in any feature column are dropped, as are rows
identical to an earlier one including the timestamp; both counts are zero on all three
datasets. Columns missing at least 30\% of the training rows are also discarded, which removes
seven channels on SWaT that are entirely absent during training, together with columns
constant on the training split and one of every pair whose absolute Pearson correlation
exceeds $0.99$, or Cram\'er's $V$ for categorical pairs. Remaining gaps are imputed with
the training median, or the mode for categorical columns.

\paragraph{Windowing.} Readings are aggregated into non-overlapping windows of a fixed
number of readings, which serve as the unit that is scored and labelled: a window is positive
if any of its readings are. Numeric features contribute their mean, standard deviation,
minimum, and maximum, while categorical features contribute their mode, entropy, distinct count, and
largest share. Windows holding fewer
than half the expected readings are dropped, which on ESA removes 85 of the 34{,}661
training windows and one window in each of the validation and test splits.
Additionally, window features are individually standardised.

\paragraph{One-class training and the two ESA labels.} The training split is rendered
one-class by removing its positive instances, ensuring that the metrics upon which detectors
calibrate are not inflated by anomalies. (Since the SWaT and HAI training splits are
already anomaly-free, no removal was required.) 
For ESA, a second variant retains these positive instances to establish a contaminated setting, 
where the detector is trained on data containing anomalies, allowing us to evaluate the masking effect caused by inflated baseline deviations. 
Furthermore, ESA annotations admit two interpretations of the positive class---wherein rare events are either treated as anomalies or grouped with nominal behaviour. 
Crossing these two labelling strategies with the two training configurations yields four distinct variants for ESA. 
Under the combined label, the validation split contains 882 positive windows across 27 segments and the test split contains 937 across 24; 
under the anomaly-only label, the validation split contains 294 positive windows across 17 segments and the test split contains 136 across 15.

\section{Methods}

\subsection{Neighbourhood graph}

Each window is a vector $x\in\mathbb{R}^{d}$ holding its $d$ features, and it receives a signature of $b$ integers, one for each of $b$ random directions. The $i$-th direction is a vector $a_i\in\mathbb{R}^{d}$ of independent standard normal entries; projecting $x$ on it gives the scalar $a_i\cdot x$; that line is cut into buckets of width $w$, shifted by a random offset $c_i$; and the integer $h_i(x)$ says which bucket $x$ falls into,
\begin{equation}
h_i(x) = \left\lfloor \frac{a_i \cdot x + c_i}{w} \right\rfloor ,
\qquad a_i \sim \mathcal{N}(0, I), \quad c_i \sim \mathcal{U}[0, w),
\end{equation}
so the signature of $x$ is $(h_1(x),\dots,h_b(x))$. 
This is the $p$-stable construction with $p=2$: the projections of two windows differ by a Gaussian whose deviation 
is their euclidean distance, so the chance of sharing a bucket falls with that distance. The bucket width is 
$w = \texttt{width\_factor} \cdot \varepsilon$, tied to the radius of interest.


Candidates are generated without exact bucket equality, by the permuted sorted-neighbourhood
scheme of Ravichandran et al.~\cite{ravichandran2005lsh} that DENCAST itself uses: the order
of the $b$ values is permuted, all windows are sorted lexicographically by the permuted
signature, and each is paired with the $\texttt{band}$ windows that follow it over several
permutations. Every
candidate pair is then verified exactly and an edge is kept when
$\lVert x_u - x_v \rVert \le \varepsilon$. The hash therefore decides only which pairs are
compared, never which are connected, and the only possible error is a near pair that no
permutation proposed.

\subsection{Clustering}

A node is a \emph{core} object when its degree reaches $\texttt{min\_pts}$. 
Clusters are the connected groups of cores, grown by max-identifier label propagation: 
a non-core node that reaches a core joins the cluster as a \emph{border} node, and a node that reaches no core remains noise.
The propagation stops when the changes fall below $\texttt{label\_change\_rate}$ times the number of directed edges, each undirected edge counting twice.

\paragraph{Calibrating the radius.} Unlike a cosine threshold, $\varepsilon$ has no universal range,
so it is set from the data, and testing quantiles of the pairwise distance distribution is one way to do it.

\subsection{Contextual subgraphs and DENCAD}

To find the contextual graph of a window to be scored (query), it is necessary to locate its closest node among the clustered ones (i.e., excluding noise points) and retrieve the corresponding cluster.
This search can be done by an exact blocked argmin using Spark: queries and pool are split into blocks, each pair of blocks is
reduced in a \texttt{pandas} task, and the per-query minima are combined.

\paragraph{Step 1: the problem.} A query is a vector $x_q \in \mathbb{R}^{d}$ and the pool
is a set $P$ of $n_p$ clustered objects, each with its own vector. The placement looks for
\begin{equation}
  \pi(q) \;=\; \arg\min_{p \in P} \lVert x_q - x_p \rVert .
\end{equation}

\paragraph{Step 2: the naive computation.} Build the full matrix $D \in
\mathbb{R}^{n_q \times n_p}$ with $D_{qp} = \lVert x_q - x_p \rVert$ and take the minimum
of each row.

\paragraph{Step 3: split the pool.} Partition $P$ into $s$ disjoint blocks,
$P = P_1 \sqcup P_2 \sqcup \dots \sqcup P_s$, and compute for each block the partial
minimum of a query,
\begin{equation}
  m_j(q) \;=\; \min_{p \in P_j} \lVert x_q - x_p \rVert .
\end{equation}

\paragraph{Step 4: Minimum combination} The minimum over a disjoint union is
the minimum of the partial minima,
\begin{equation}
  \min_{p \in P} \lVert x_q - x_p \rVert \;=\; \min_{j = 1, \dots, s} m_j(q) ,
  \label{eq:minsplit}
\end{equation}

\paragraph{Step 5: split the queries as well.} With $Q = Q_1 \sqcup \dots \sqcup Q_r$ there
are $r \times s$ block pairs $(Q_i, P_j)$, and each is independent of the others.

\paragraph{Step 6: map and reduce.} Each of the $r \cdot s$ pairs is computed on its own
and emits $m_j(q)$ for every $q \in Q_i$; the results are grouped by $q$ and reduced by
taking the minimum of its $s$ values.

Its contextual subgraph is induced from the cluster of the identified nearest node. 
Specifically, the query is attached to both the nearest object and any cluster members 
within distance $\varepsilon$. The subgraph is then formed by extending this attachment 
set along the cluster's existing edges up to an $(L-1)$-hop neighbourhood, which gives 
$L$ graph convolutions their complete receptive field. The final graph comprises 
these nodes, the cluster's edges among them, the query's attachment edges, and the 
query itself designated as the central target node.

The encoder is a graph convolution with the symmetric normalisation of~\cite{kipf2017gcn},
\begin{equation}
  H^{(l+1)} = \sigma\!\left( \hat{A} H^{(l)} W^{(l)} \right), \qquad
  \hat{A} = \tilde{D}^{-1/2}\tilde{A}\tilde{D}^{-1/2},
\end{equation}
and the decoder reconstructs the centre's own attributes. The score of a window is the
distance between the centre and its reconstruction.
\begin{equation}
  s(x) = \sqrt{\tfrac{1}{d}\lVert \hat{x} - x \rVert^{2}} ,
\end{equation}

\subsection{Training examples}

Training examples are constructed as follows. Each training window—including noise points—is 
treated as a query and excluded from its own context. It is attached to the contextual 
subgraph using the cluster of its nearest non-noise object, ensuring that training examples 
are generated identically to those at inference time.

\subsection{Distributed fit}

Training examples are sharded by cluster, so a replica owns whole clusters and never needs
a member it does not hold. Four
replicas then fit the same network with \texttt{DistributedDataParallel} over gloo and
gradients are averaged at every step.

\subsection{Baselines}

The \emph{plain autoencoder} is DENCAD with $\hat{A} = I$: same architecture,
hyperparameters, and score, fitted on single windows, which isolates the contribution of
the graph. It is useful to determine if the neighbourhood graph structure is really providing 
value to the anomaly detection task. \emph{K-Means} scores the distance to the nearest centroid, 
with $k$ selected via validation. \emph{DENCAST distance} uses density-based clustering but 
omits the autoencoder, scoring the distance to the closest clustered object. The \emph{graph
autoencoder} is DENCAD without the clustering: it applies the same network over the complete 
$L$-hop neighbourhood of the window, thereby matching the exact output that a full-graph 
convolution would produce at the centre.

\paragraph{Comparison} We present threshold-free measures as our primary comparison, 
alongside window-level precision and recall, while reporting the number of 
covered anomalous episodes as a supplementary information.


\subsection{Evaluation}

All five anomaly detectors are thresholded by the same technique that uses the sigma rule. It is divied in three phases:
\begin{equation}
  z(x) = \frac{s(x) - \mu}{\sigma}, \qquad \text{flag when } z(x) > k .
\end{equation}

In phase one, the model is fitted on the training set, $\mu$ and $\sigma$ are measured on its
training scores, and $k$ is chosen among ten candidates
between $0.5$ and $4$ via validation $F_1$; for the two autoencoders, the epoch count is
selected in the same phase using early stopping on validation average precision, with a
patience of $20$ and a maximum of $200$ epochs. In phase two, the model is refitted on the 
combined training and validation windows for that epoch count (normal or contaminated depending on the specific experiment), and $\mu$ and $\sigma$ are
recomputed on this larger set. Finally, in phase three, the test set is used to compute the final measures for the comparison.

The hyperparameter grids vary significantly in size, requiring different search strategies. 
K-Means has a single hyperparameter, the number of centroids $k \in \{2, 4, 8, 16, 32, 64\}$, and 
all six values are evaluated; the plain autoencoder comprises sixteen combinations and is 
likewise tested exhaustively. The remaining grids (DENCAST distance, DENCAD, and Graph Autoencoder) 
are too large; thus, due to computational constraints, fifteen configurations were sampled 
at random from each. In every case, the final configuration is selected based on validation average precision.

The searched axes are $b \in \{16, 24\}$, permutations $\in \{10, 20\}$, band
$\in \{20, 50\}$, $w/\varepsilon \in \{1, 2\}$, $\texttt{min\_pts} \in \{10, 5\}$ for the
graph; hidden $\in \{64, 128\}$, latent $\in \{16, 32\}$, layers $\in \{2, 3\}$ and
learning rate $\in \{10^{-3}, 3\cdot 10^{-3}\}$ for the autoencoder. Batch size 128, patience 20, a ceiling
of 200 epochs, four replicas, Adam and one seed are fixed.

\paragraph{Metrics.} Average precision and ROC-AUC are reported as the primary comparison
because they are threshold-free. Precision, recall, and $F_1$ are evaluated at the
\emph{window} level, while the number of anomalous segments covered at
least once is reported separately.

\paragraph{Implementation.} PySpark 3.5.9 on a local four-worker session with the PyTorch 
gloo backend, hosted on a single machine with four cores and 16\,GB of memory. 
The pipeline comprises 66 DVC stages.

\section{Results}

\begin{table}[t]
  \centering
  \small
  \caption{Test average precision, best per row in bold. Five detectors, six processed
  datasets, one seed.}
  \begin{tabular}{lrrrrr}
    \toprule
    & DENCAD & graph AE & plain AE & DENCAST $\sigma$ & K-Means \\
    \midrule
SWaT & 0.8484 & 0.8499 & \textbf{0.8545} & 0.8197 & 0.8340 \\
HAI & 0.5734 & 0.5786 & \textbf{0.5972} & 0.5245 & 0.5532 \\
ESA rare & 0.2313 & 0.2307 & \textbf{0.2424} & 0.2280 & 0.2255 \\
ESA rare, cont. & 0.1708 & \textbf{0.1907} & 0.1671 & 0.1681 & 0.1623 \\
ESA & \textbf{0.0533} & 0.0384 & 0.0409 & 0.0390 & 0.0315 \\
ESA, cont. & \textbf{0.0263} & 0.0245 & 0.0253 & 0.0247 & 0.0246 \\
    \bottomrule
  \end{tabular}
  \label{tab:ap}
\end{table}

\begin{table}[t]
  \centering
  \small
  \caption{Test ROC-AUC, best per row in bold. Five detectors, six processed datasets,
  one seed.}
  \begin{tabular}{lrrrrr}
    \toprule
    & DENCAD & graph AE & plain AE & DENCAST $\sigma$ & K-Means \\
    \midrule
SWaT & 0.9239 & 0.9257 & \textbf{0.9277} & 0.9066 & 0.9094 \\
HAI & 0.8393 & \textbf{0.8525} & 0.8406 & 0.8027 & 0.8174 \\
ESA rare & 0.6515 & 0.6510 & \textbf{0.6670} & 0.6543 & 0.6449 \\
ESA rare, cont. & 0.6341 & \textbf{0.6519} & 0.6405 & 0.6270 & 0.6098 \\
ESA & 0.6109 & 0.6132 & 0.5932 & \textbf{0.6290} & 0.6066 \\
ESA, cont. & 0.5965 & 0.6041 & \textbf{0.6063} & 0.5964 & 0.5936 \\
    \bottomrule
  \end{tabular}
  \label{tab:roc}
\end{table}

Overall, autoencoder-based detectors slightly outperform distance-based methods (K-Means and DENCAST distance) 
across nearly all datasets under both metrics, while little to no performance difference is observed between DENCAST distance and K-Means.

Among the autoencoders, no clear overall winner emerges, and incorporating the graph does not 
always pay for itself: the plain autoencoder, which operates on single observations without 
graph structure, achieves the highest score on three of the six datasets under either metric. 
DENCAD achieves the highest average precision on both ESA variants that consider rare events 
as legitimate---where all detectors yield scores between $0.025$ and $0.053$---and its margin 
is substantially wider than indicated by ROC-AUC.

\section{Conclusion}

For anomaly detection over multivariate sensor telemetry from industrial plants and spacecraft, 
reconstruction-based models prove to be slightly better than distance-based methods. 

No clear performance difference was observed between K-Means and DENCAST distance. 
Although DENCAST theoretically yields superior clustering, the distance to an individual cluster object might be less informative than the distance to a centroid; however, these explanations remain speculative.

The construction of a neighbourhood graph does not consistently justify its computational overhead, 
and its overall utility appears to be highly dataset-dependent. 
In this study, DENCAD does not exhibit a performance advantage that offsets the cost of graph building.
However, this finding specifically pertains to graphs constructed via locality-sensitive hashing (LSH) in feature space; 
evaluating DENCAD on inherently graph-structured datasets remains an important avenue for future work.

One limitation bounds the comparison: the hyperparameter grids are explored unequally. 
While the plain autoencoder and K-Means algorithms exhaustively enumerate all their combinations,
DENCAD samples fifteen at random due to the severe combinatorial explosion of its hyperparameter space.
Consequently, DENCAD's current figures represent a lower bound on its potential performance, 
leaving open the possibility that a fully searched grid could close or reverse the gap.
Although conducting an exhaustive search is computationally expensive and highlights DENCAD's vulnerability to hyperparameter explosion, 
evaluating the complete grid remains a worthwhile endeavor for future work.


All five detectors compute an \emph{unweighted} Euclidean distance within the same feature space—whether to a centroid, 
a core object, or between a window and its reconstruction—meaning that while the context and model architecture varied, 
the aggregation method remained constant. Consequently, every feature contributes equally to the distance calculation.
Scaling each feature's error by its typical magnitude during training is identified as an important modification for future work, 
and all five detectors should be re-evaluated under this weighted formulation.

\bibliographystyle{unsrt}

{\footnotesize
\begin{thebibliography}{10}
\scriptsize
\setlength{\itemsep}{0pt}
\setlength{\parskip}{0pt}
\setlength{\parsep}{0pt}
\bibitem{corizzo2019dencast}
Roberto Corizzo, Gianvito Pio, Michelangelo Ceci, and Donato Malerba.
\newblock {DENCAST}: distributed density-based clustering for multi-target regression.
\newblock {\em Journal of Big Data}, 6:43, 2019.
\bibitem{ester1996dbscan}
Martin Ester et~al.
\newblock A density-based algorithm for discovering clusters in large spatial databases with noise.
\newblock In {\em KDD}, pages 226--231, 1996.
\bibitem{charikar2002simhash}
Moses S. Charikar.
\newblock Similarity estimation techniques from rounding algorithms.
\newblock In {\em STOC}, pages 380--388, 2002.
\bibitem{datar2004pstable}
Mayur Datar, Nicole Immorlica, Piotr Indyk, and Vahab S. Mirrokni.
\newblock Locality-sensitive hashing scheme based on $p$-stable distributions.
\newblock In {\em SoCG}, pages 253--262, 2004.
\bibitem{ravichandran2005lsh}
Deepak Ravichandran, Patrick Pantel, and Eduard Hovy.
\newblock Randomized algorithms and {NLP}: using locality sensitive hash functions for high speed noun clustering.
\newblock In {\em ACL}, pages 622--629, 2005.
\bibitem{ding2019dominant}
Kaize Ding et~al.
\newblock Deep anomaly detection on attributed networks.
\newblock In {\em SDM}, pages 594--602, 2019.
\bibitem{kipf2017gcn}
Thomas N. Kipf and Max Welling.
\newblock Semi-supervised classification with graph convolutional networks.
\newblock In {\em ICLR}, 2017.
\bibitem{hamilton2017graphsage}
William L. Hamilton, Rex Ying, and Jure Leskovec.
\newblock Inductive representation learning on large graphs.
\newblock In {\em NeurIPS}, pages 1024--1034, 2017.
\bibitem{chiang2019clustergcn}
Wei-Lin Chiang et~al.
\newblock {Cluster-GCN}: an efficient algorithm for training deep and large graph convolutional networks.
\newblock In {\em KDD}, pages 257--266, 2019.
\newblock Towards a rigorous evaluation of time-series anomaly detection.
\newblock In {\em AAAI}, pages 7194--7201, 2022.
\newblock Current time series anomaly detection benchmarks are flawed and are creating the illusion of progress.
\newblock {\em IEEE Transactions on Knowledge and Data Engineering}, 2023.
\bibitem{mathur2016swat}
Aditya P. Mathur and Nils Ole Tippenhauer.
\newblock {SWaT}: a water treatment testbed for research and training on {ICS} security.
\newblock In {\em CySWater}, pages 31--36, 2016.
\bibitem{shin2020hai}
Hyeok-Ki Shin et~al.
\newblock {HAI} 1.0: {HIL}-based augmented {ICS} security dataset.
\newblock In {\em USENIX CSET}, 2020.
\bibitem{kotowski2024esa}
Krzysztof Kotowski et~al.
\newblock European {Space} {Agency} benchmark for anomaly detection in satellite telemetry.
\newblock {\em arXiv preprint arXiv:2406.17826}, 2024.
\newblock Accurate, large minibatch {SGD}: training {ImageNet} in 1 hour.
\newblock {\em arXiv preprint arXiv:1706.02677}, 2017.
\end{thebibliography}
}

\end{document}
